% !TeX program = xelatex
\documentclass{UAmathtalk}
\usepackage{setspace}
\renewcommand{\when}{12:00 p.m.}
\renewcommand{\where}{Hudson 0110}
\renewcommand{\where}{Hudson 0110.\\ \href{https://albany.zoom.us/j/81219478289?pwd=N5fy7vxKBXqlQ4mAq4KaeaAM5iAHi4.1}{Zoom ID: 812\,1947\,8289 - Passcode: 233444}}
\author{Gonzalo Ortega Carpintero}
\address{City St Georges, University of London}
\urladdr{https://gonzaloortega.com/}
\title{Persistent fundamental groups and lower-central-quotient barcodes}
\date{Friday, November 13, 2026}


\begin{document}

\maketitle
%\renewcommand*{\abstractsize}{\normalsize}
\renewcommand*{\abstractsize}{\fontsize{13.5}{16.2}\selectfont}
\begin{abstract}
    {
        We present an efficient algorithm for computing the persistent fundamental group of a filtered CW-complex in terms of group presentations. From these presentations, we derive persistent invariants associated with the lower central series and obtain a family of barcodes that capture topological features not detected by persistent homology. This framework extends the scope of topological data analysis by enabling the study of non-abelian aspects of the fundamental group from filtered data.    }
\end{abstract}
\end{document}



